% AWHDL チュートリアル（日本語）
% ビルド: latexmk -lualatex awhdl_tutorial_ja.tex
\documentclass[a4paper,11pt]{ltjsarticle}

\usepackage{luatexja-fontspec}
\usepackage{amsmath}
\usepackage{xcolor}
\usepackage{listings}
\usepackage{booktabs}
\usepackage{longtable}
\usepackage{tabularx}
\usepackage{array}
\usepackage[most]{tcolorbox}
\usepackage{tikz}
\usetikzlibrary{arrows.meta,positioning,fit,shapes.geometric}
\usepackage[hidelinks,pdfusetitle]{hyperref}
\usepackage{geometry}
\geometry{margin=24mm}

\definecolor{kw}{RGB}{24,94,153}
\definecolor{cm}{RGB}{94,109,126}
\definecolor{st}{RGB}{39,137,126}
\definecolor{bg}{RGB}{246,248,250}
\definecolor{note}{RGB}{226,239,249}
\definecolor{warn}{RGB}{251,239,222}

\lstdefinelanguage{AWHDL}{
  morekeywords={entity,architecture,is,of,port,in,out,inout,end,begin,
    device,generic,signal,timer,period,budget,barrier,process,timeout,on,
    if,then,elsif,else,parallel,assert,always,never,null,and,or,not,true,false},
  sensitive=true,
  morecomment=[l]{--},
  morestring=[b]",
}
\lstset{
  basicstyle=\ttfamily\small,
  keywordstyle=\color{kw}\bfseries,
  commentstyle=\color{cm}\itshape,
  stringstyle=\color{st},
  backgroundcolor=\color{bg},
  frame=single,
  rulecolor=\color{black!15},
  numbers=left,
  numberstyle=\tiny\color{cm},
  xleftmargin=2em,
  framexleftmargin=1.5em,
  columns=fullflexible,
  keepspaces=true,
  breaklines=true,
  showstringspaces=false,
  tabsize=4,
}
\lstdefinestyle{shell}{language={},numbers=none,xleftmargin=0.5em,framexleftmargin=0em,
  keywordstyle={},commentstyle={}}

\newtcolorbox{point}[1][]{colback=note,colframe=kw!60,fonttitle=\bfseries,title=#1,breakable}
\newtcolorbox{caution}[1][]{colback=warn,colframe=orange!70!black,fonttitle=\bfseries,title=#1,breakable}

\newcommand{\ex}[1]{\texttt{examples/tutorial/#1}}

\title{AWHDL チュートリアル\\\large Agentic Workflow Harness Description Language v0.1}
\author{AI Conductor プロジェクト}
\date{2026年9月24日}

\begin{document}
\maketitle

\begin{abstract}
AWHDL は、複数の AI・MCP ツール・決定的プログラム・人間を組み合わせた反復ループを、
型付き・イベント駆動の設計として記述する言語です。本チュートリアルでは、
実機で動作を確認した 6 つの設計例を順にたどり、\texttt{entity}／\texttt{architecture}、
device と route の束縛、signal と delta cycle、タイマーと予算、フィードバックループ、
並列実行とバリア、情報フロー検査までを学びます。
\end{abstract}

\tableofcontents
\clearpage

%======================================================================
\section{はじめに}

\subsection{この文書の位置付け}

本書は非規範の入門書です。v0.1 の実装範囲は \texttt{docs/PROFILE\_v0.1\_ja.md} だけが定め、
言語・実行・安全の規則は LANGUAGE\_SPEC、RUNTIME\_SPEC、SECURITY\_SPEC、IR\_SPEC が定めます。
本書の記述がそれらと食い違う場合は、規範文書が優先します。

本書の例はすべて \ex{} にあり、2026年9月24日に次の 2 つで確認しました。
\begin{itemize}
  \item \texttt{aic check} による構文・構造・情報フローの静的検査
  \item \texttt{aiconductor run-design} による実機での実行（27B モデル、Prolog、Lean、KDB を使用）
\end{itemize}

\subsection{AWHDL が解く問題}

AI にツールを使わせる仕組みでは、「どの AI がどのデータを受け取るか」「いつ再試行し、いつ止めるか」
「何をもって完了とするか」がプロンプトの中に埋もれがちです。AWHDL はこれらを設計として書き出し、
実行前に検査できるようにします。

\begin{point}[AWHDL の基本的な考え方]
\begin{itemize}
  \item AI やツールは \textbf{device}（入出力を持つ部品）として扱う。
  \item データや状態は \textbf{signal} が運ぶ。
  \item \textbf{process} はイベントに反応して動く。
  \item 実行を続けてよいか、外部へ送ってよいか、完了したかは、AI ではなく
        \textbf{AI Conductor ランタイム}が宣言された規則に従って決める。
\end{itemize}
\end{point}

構文は VHDL に似ていますが、AWHDL はハードウェアを記述する言語ではありません。
名前の H は Harness（複数の実行主体を一つの制御構造に束ねるもの）を表します。

\subsection{全体の構成}

\begin{center}
\begin{tikzpicture}[
  box/.style={draw,rounded corners,fill=bg,minimum width=30mm,minimum height=9mm,align=center,font=\small},
  arr/.style={-{Stealth[length=2mm]},thick}]
\node[box] (src) {AWHDL 設計\\\texttt{.awhdl}};
\node[box,right=10mm of src] (chk) {\texttt{aic check}\\静的検査};
\node[box,right=10mm of chk] (rt) {\texttt{run-design}\\束縛と delta cycle};
\node[box,below=12mm of rt] (eng) {AI Conductor エンジン\\権限・Effect・承認・監査};
\node[box,left=10mm of eng] (dev) {device\\LLM / MCP};
\node[box,right=10mm of rt] (out) {出力 JSON};
\draw[arr] (src) -- (chk);
\draw[arr] (chk) -- (rt);
\draw[arr] (rt) -- (out);
\draw[arr] (rt) -- (eng);
\draw[arr] (eng) -- (dev);
\end{tikzpicture}
\end{center}

\texttt{run-design} は設計を検査・束縛したうえで delta cycle として実行します。
device への呼び出しは、通常の実行と同じエンジンを通るため、権限（capability）の検査、
Effect の記録、承認、監査ログがすべて適用されます。

%======================================================================
\section{準備}

\subsection{必要なもの}

\begin{itemize}
  \item Rust（2024 edition）と Cargo
  \item \texttt{awhdl} リポジトリ（本書の例とランタイムを含む）
  \item 実行（\texttt{run-design}）には、AI Conductor のプロジェクトルートと、
        そこで設定された route（モデルや MCP サーバー）
\end{itemize}

静的検査（\texttt{aic check}）だけなら、モデルや MCP サーバーは不要です。

\subsection{ビルド}

\begin{lstlisting}[style=shell]
cd awhdl
cargo build -p conductor-cli -p aiconductor
# aic         : 構文と静的検査
# aiconductor : AI Conductor ランタイム（run-design を含む）
\end{lstlisting}

AI Conductor のワークスペースでは、ビルド成果物を OneDrive の外に置くため、次の設定を使います。

\begin{lstlisting}[style=shell]
export CARGO_TARGET_DIR=~/.cache/aiconductor/cargo-target
\end{lstlisting}

以下では実行ファイルを \texttt{aic}、\texttt{aiconductor} と略記します。

\subsection{検査と実行のコマンド}

\begin{lstlisting}[style=shell]
# 静的検査（終了状態 0 = 合格）
aic check examples/tutorial/01_hello.awhdl
aic check examples/tutorial/01_hello.awhdl --output json

# 実行（--project-root は AI Conductor のルート）
aiconductor --project-root <AIConductor> run-design \
    awhdl/examples/tutorial/01_hello.awhdl --input 'task=...'
\end{lstlisting}

\texttt{--input 名前=値} は入力ポートに値を与えます。値が JSON として解釈できれば JSON、
そうでなければ文字列になります。複数のポートには \texttt{--input} を繰り返します。

\begin{table}[h]
\centering
\caption{\texttt{aic check} の終了状態}
\begin{tabular}{cl}
\toprule
状態 & 意味 \\
\midrule
0 & 実装済みの検査に合格 \\
1 & ソースファイルを読み取れない \\
2 & 構文解析に失敗 \\
3 & 構造または情報フローの検査に失敗 \\
\bottomrule
\end{tabular}
\end{table}

%======================================================================
\section{最初の設計：entity と architecture}

\subsection{設計}

\lstinputlisting[language=AWHDL,caption={\ex{01\_hello.awhdl}}]{../../examples/tutorial/01_hello.awhdl}

\begin{description}
  \item[\texttt{entity}] 設計の外部インターフェースです。\texttt{port} に入力（\texttt{in}）と
    出力（\texttt{out}）を並べます。
  \item[型 \texttt{text<internal>}] 型名の後の \texttt{<\dots>} は\textbf{機密区分}です。
    \texttt{public}、\texttt{internal}、\texttt{restricted} の 3 段階があり、第\ref{sec:flow}節の
    情報フロー検査に使われます。
  \item[\texttt{architecture}] entity の中身です。\texttt{is} と \texttt{begin} の間に宣言を、
    \texttt{begin} と \texttt{end} の間に process を書きます。
  \item[\texttt{device}] AI やツールを表す部品です。種類は \texttt{agent}、\texttt{mcp}、
    \texttt{deterministic} の 3 つです。\texttt{generic} で、どの \textbf{route}
    （AI Conductor の設定にある実行経路）に結び付けるか、どこで動くか（\texttt{location}）を宣言します。
  \item[\texttt{process(task)}] 括弧内は\textbf{感度リスト}です。\texttt{task} が変化すると本体が実行されます。
  \item[\texttt{coder.run(task) -> answer;}] device を呼び出し、結果を \texttt{answer} に書き込みます。
    agent device の場合、メソッド名（ここでは \texttt{run}）は結果に影響しません。
\end{description}

\subsection{実行}

\begin{lstlisting}[style=shell]
aiconductor --project-root . run-design awhdl/examples/tutorial/01_hello.awhdl \
  --input 'task=Rustで2つのi32を安全に足す関数を1つだけ書いてください。説明は1行で。'
\end{lstlisting}

\begin{lstlisting}[style=shell,caption={実行結果（27B モデル、約 4 秒）}]
{
  "status": "quiescent",
  "reason": null,
  "outputs": {
    "answer": "```rust\nfn safe_add(a: i32, b: i32) -> Option<i32> {\n    a.checked_add(b)\n}\n```\n\n`checked_add` はオーバーフロー時に `None` を返すため、..."
  },
  "delta_cycles": 2,
  "device_calls": 1
}
\end{lstlisting}

\begin{description}
  \item[\texttt{status}] \texttt{quiescent} は、処理すべきイベントがなくなって自然に終わったことを表します。
    予算で打ち切られた場合は \texttt{exhausted} になり、\texttt{reason} に理由が入ります。
  \item[\texttt{outputs}] 出力ポートの最終値です。
  \item[\texttt{delta\_cycles}] 実行した delta cycle の数です（第\ref{sec:delta}節）。
  \item[\texttt{device\_calls}] device を呼び出した回数です。
\end{description}

%======================================================================
\section{決定的な判定：deterministic device と if}

モデルの出力をそのまま「合格」とみなすと、モデルの誤りがそのまま結論になります。
AWHDL では、ツールが返した構造化データを使って、判定を設計の中で決定的に行えます。

\lstinputlisting[language=AWHDL,caption={\ex{02\_prolog\_verdict.awhdl}}]{../../examples/tutorial/02_prolog_verdict.awhdl}

\begin{description}
  \item[\texttt{deterministic}] 読み取り専用の MCP route にだけ結び付けられる device です。
    書き込みを伴う route（MATLAB や KDB など）には \texttt{mcp} を使います。
  \item[\texttt{signal result : data<internal>;}] 設計内部の信号です。MCP device の結果は
    JSON オブジェクトで、\texttt{result.text}（ツールの出力文字列）や
    \texttt{result.structured}（ツールが返した JSON）としてフィールドを読めます。
  \item[\texttt{timeout 1 min}] 呼び出しの制限時間です。時間の単位は
    \texttt{ms}、\texttt{sec}、\texttt{min}、\texttt{hour}、\texttt{day} です。
  \item[\texttt{if \dots\ elsif \dots\ else \dots\ end if;}] 条件分岐です。
  \item[\texttt{rules.failed}] 呼び出しが失敗したときに起きるイベントです。成功時は
    \texttt{rules.done}、制限時間の超過時は \texttt{rules.timeout} と \texttt{rules.failed} が起きます。
\end{description}

\subsection{入力の渡し方}

v0.1 の式にはオブジェクトを組み立てる構文がないため、device に渡すのは主に 1 つの文字列です。
ツールの引数（Prolog のプログラムと問い合わせなど）は、ランタイムのアダプタが入力文字列から取り出します。
Prolog では、バッククォートで囲んだ節の並びをプログラム、最後の単一のゴールを問い合わせとして扱います。

\begin{lstlisting}[style=shell]
aiconductor --project-root . run-design awhdl/examples/tutorial/02_prolog_verdict.awhdl \
  --input 'task=`parent(tom,bob). parent(bob,ann). grandparent(X,Z) :- parent(X,Y), parent(Y,Z).` `grandparent(tom,ann).`'
# => "verdict": "holds"          （delta_cycles 3, device_calls 1）

aiconductor --project-root . run-design awhdl/examples/tutorial/02_prolog_verdict.awhdl \
  --input 'task=`parent(tom,bob). grandparent(X,Z) :- parent(X,Y), parent(Y,Z).` `grandparent(tom,ann).`'
# => "verdict": "does not hold"
\end{lstlisting}

\begin{point}[ポイント]
「成り立つかどうか」は Prolog が返した \texttt{solution\_count} で決まり、モデルは関与しません。
モデルに判断させるのは、決定的に判定できない部分だけにするのが AWHDL での基本的な設計です。
\end{point}

%======================================================================
\section{delta cycle、タイマー、予算}\label{sec:delta}

\subsection{delta cycle}

AWHDL の実行は \textbf{delta cycle} の繰り返しです。1 つの delta cycle では次の順に処理します。

\begin{enumerate}
  \item 直前に起きたイベントに感度を持つ process を選ぶ。
  \item 選んだ process をすべて、同じ時点の値のスナップショットに対して実行する。
        process 内の代入（\texttt{<=}）や device の結果は、すぐには反映されない。
  \item delta cycle の終わりに、書き込みをまとめて反映する。値が\textbf{実際に変わった} signal だけが
        次のイベント（\texttt{名前} と \texttt{名前.changed}）を起こす。
  \item 同じ delta cycle で 2 つの process が同じ signal に異なる値を書くとエラーになる。
  \item \texttt{assert always}／\texttt{assert never} を評価する。
\end{enumerate}

イベントが残らず、タイマーもなければ、実行は \texttt{quiescent} で終わります。

\subsection{タイマーと予算}

\lstinputlisting[language=AWHDL,caption={\ex{03\_timer\_budget.awhdl}}]{../../examples/tutorial/03_timer_budget.awhdl}

\begin{description}
  \item[\texttt{timer tick : period 1 sec;}] 1 秒ごとに \texttt{tick} イベントを起こします。
    タイマーを使う設計には、\texttt{iterations} か \texttt{wall\_time} の予算が必須です。
  \item[\texttt{budget}] 実行の上限です。v0.1 で使える項目は \texttt{iterations}（delta cycle 数）、
    \texttt{wall\_time}（経過時間）、\texttt{tool\_calls}、\texttt{model\_calls} です。
  \item[\texttt{signal count : data<internal> := 0;}] 初期値付きの signal です。初期値は定数に限ります。
  \item[\texttt{assert always (count <= 10);}] すべての delta cycle の終わりに成り立つべき条件です。
    破られると実行は失敗します。
\end{description}

\begin{lstlisting}[style=shell,caption={実行結果（デバイスなし、約 3 秒）}]
{ "status": "exhausted", "reason": "wall_time budget",
  "outputs": { "ticks": 2 }, "delta_cycles": 6, "device_calls": 0 }
\end{lstlisting}

\begin{caution}[設計の予算とランタイムの予算]
設計の \texttt{budget} は、AI Conductor の \texttt{config/runtime.toml}（\texttt{[loop]} 節）の上限を
さらに狭めるだけで、広げることはできません。例えばランタイムの \texttt{max\_iterations} が 8 のとき、
設計で \texttt{iterations <= 20} と書いても 8 delta cycle で打ち切られます。
タイマー間隔を 500\,ms にすると、この例は設計の \texttt{wall\_time} より先に
ランタイムの上限（\texttt{iteration budget exhausted}）で止まります。
\end{caution}

%======================================================================
\section{フィードバックループ}

検証して、結果によって再実行する。AWHDL では、このループをイベントの連鎖として書きます。

\lstinputlisting[language=AWHDL,caption={\ex{04\_feedback\_loop.awhdl}}]{../../examples/tutorial/04_feedback_loop.awhdl}

\begin{enumerate}
  \item \texttt{code} が与えられると \texttt{round <= 1}。
  \item \texttt{round} が変わるたびに、条件を満たせば KDB/q を呼ぶ。
  \item 呼び出しが成功すると \texttt{q.done} が起き、\texttt{round} を 1 増やす。
  \item \texttt{round} が 4 になると呼び出さなくなり、イベントが尽きて終わる。
\end{enumerate}

\begin{caution}[signal のイベントと device のイベント]
signal のイベントは値が\textbf{変わったとき}だけ起きます。同じ計算を繰り返すと結果が変わらないため、
\texttt{process(result)} ではループが 1 回で止まってしまいます。呼び出しの完了ごとに必ず起きる
\texttt{q.done} を使うのはこのためです。
\end{caution}

\begin{lstlisting}[style=shell]
aiconductor --project-root . run-design awhdl/examples/tutorial/04_feedback_loop.awhdl \
  --input 'code=sum til 101'
# => "rounds": 3, "last": "{ \"ok\": true, ... \"stdout\": \"5050\\n\" ... }"
#    status quiescent, delta_cycles 8, device_calls 3
\end{lstlisting}

\texttt{tool\_calls <= 5} は、万一条件を書き誤っても呼び出しが 5 回を超えないための安全策です。
\texttt{assert never (round > 4);} は「このループは 4 回目以降に進まない」という設計者の意図を、
実行時に検査できる形で残します。

%======================================================================
\section{並列実行とバリア}

独立した検査は同時に始め、両方がそろった時点で次へ進みたいことがあります。

\lstinputlisting[language=AWHDL,caption={\ex{05\_parallel\_barrier.awhdl}}]{../../examples/tutorial/05_parallel_barrier.awhdl}

\begin{description}
  \item[\texttt{parallel \dots\ end parallel;}] 中の呼び出しは論理的に並列で、結果は同じ delta cycle の
    終わりにまとめて反映されます（v0.1 のランタイムは実際には 1 つずつ順に呼び出します）。
    2 つの分岐が同じ signal に書き込むことはできません（E215）。
  \item[\texttt{barrier both (rule\_result, proof\_result);}] 列挙した signal がすべて書き込まれると
    \texttt{both.ready} が起きます。
\end{description}

\begin{lstlisting}[style=shell]
aiconductor --project-root . run-design awhdl/examples/tutorial/05_parallel_barrier.awhdl \
  --input 'rules_task=`edge(a,b). path(X,Y) :- edge(X,Y).` `path(a,b).`' \
  --input 'proof_task=theorem two : 1 + 1 = 2 := by rfl'
# => "summary": "rules and proof both returned"   （device_calls 2）

# 誤った定理を与えると Lean が拒否し、proof_result は書き込まれない
  --input 'proof_task=theorem bad : 2 + 2 = 5 := by rfl'
# => "summary": "the proof was rejected"
\end{lstlisting}

失敗した呼び出しは出力 signal に書き込まないため、バリアは成立しません。
失敗は \texttt{prover.failed} で別に扱います。

%======================================================================
\section{情報フローと安全}\label{sec:flow}

\subsection{機密区分と配置}

signal とポートの型には機密区分（\texttt{public} < \texttt{internal} < \texttt{restricted}）を付けます。
区分のない型は \texttt{restricted} として扱われます。device は \texttt{location}
（\texttt{local} または \texttt{cloud}）を宣言し、route の実際の配置と一致しなければ束縛に失敗します。

\lstinputlisting[language=AWHDL,caption={\ex{06\_information\_flow.awhdl}}]{../../examples/tutorial/06_information_flow.awhdl}

\begin{lstlisting}[style=shell]
$ aic check awhdl/examples/tutorial/06_information_flow.awhdl
awhdl/examples/tutorial/06_information_flow.awhdl:14:9: AWHDL-E301: Restricted data cannot flow to cloud device reviewer
$ echo $?
3
\end{lstlisting}

\texttt{restricted} のデータがクラウドの device に流れる経路は、実行前に拒否されます。
ランタイムも呼び出しの直前に同じ検査を行うため、静的検査をすり抜けた場合でも送信されません。

\subsection{さらに厳しくする}

\begin{itemize}
  \item \texttt{assert never (internal -> cloud);} を書くと、\texttt{internal} 以上のデータも
        クラウドへ送れなくなります。
  \item device に \texttt{clearance => public} などを付けると、その device が受け取れる区分の上限を
        決められます（超えると E302）。
  \item 低い区分の signal に高い区分のデータを代入することはできません（E303）。
        v0.1 には区分を下げる仕組み（declassifier）はありません。
\end{itemize}

\subsection{設計の外にある安全機構}

次の規則は設計ではなく、AI Conductor の設定とランタイムが強制します。
\begin{itemize}
  \item \textbf{capability}：どの route がどのツール・ファイル・モデルを使えるか。
        設定は \nolinkurl{config/capabilities.toml} にあります。
  \item \textbf{Effect と承認}：書き込みを伴う呼び出しの記録と、外部への書き込み前の人間の承認。
  \item \textbf{監査}：実行ごとの \texttt{.runtime/runs/<id>/events.jsonl}。
\end{itemize}
設計から device を呼び出すと、これらがすべて適用されます。
設計の側からこれらを無効にする構文はありません。

%======================================================================
\section{device と route}

\subsection{束縛の規則}

\begin{table}[h]
\centering\small
\caption{device の種類と結び付けられる route}
\begin{tabular}{lll}
\toprule
device の種類 & 結び付けられる route & 例 \\
\midrule
\texttt{agent} & モデルの route、または Codex & \texttt{initial\_coding}、\texttt{deep\_reasoning} \\
\texttt{mcp} & MCP の route（Codex 以外） & \texttt{calculation\_graphing}、\texttt{table\_analysis} \\
\texttt{deterministic} & 読み取り専用の MCP の route & \texttt{logic\_rules}、\texttt{logic\_proof}、\texttt{text\_filtering} \\
\bottomrule
\end{tabular}
\end{table}

主制御（planner）の route（\texttt{loop}）には結び付けられません。

\subsection{AI Conductor の route（2026年9月時点）}

{\small
\begin{longtable}{p{44mm}p{27mm}p{72mm}}
\toprule
route & 種類 & 実行先 \\
\midrule
\endhead
\texttt{initial\_coding} & モデル & 192.168.6.10 の RTX 4090（Qwen3.8 27B） \\
\texttt{japanese\_polishing} & モデル & 192.168.6.10（llm-jp-4 8B） \\
\texttt{vision} & モデル & 192.168.6.20 の RX 9060 XT（Qwen3-VL 8B） \\
\texttt{logic\_rules} & MCP（読み取り） & 192.168.6.10 の SWI-Prolog \\
\texttt{logic\_proof} & MCP（読み取り） & 192.168.6.10 の Lean 4.33.1 \\
\texttt{text\_filtering} & MCP（読み取り） & ローカルの制限付き Filter \\
\texttt{table\_analysis} & MCP（書き込み） & 192.168.6.10 の KDB/q \\
\texttt{calculation\_graphing} & MCP（書き込み） & 192.168.6.10 の MATLAB R2026a \\
\texttt{deep\_reasoning} & Codex（cloud） & Codex CLI（読み取り専用） \\
\texttt{open\_data\_acquisition} & Codex（cloud） & Codex CLI（承認が必要な外部書き込み） \\
\bottomrule
\end{longtable}}

\subsection{呼び出しと引数}

\begin{itemize}
  \item 引数が 1 つの文字列なら、それが device への入力になります。
  \item MCP の route では、メソッド名がその route のツール名と一致すれば、そのツールを使います
        （例：\texttt{rules.run\_prolog}、\texttt{prover.check\_lean\_code}、\texttt{q.run\_q}）。
  \item ツールの引数はアダプタが入力から取り出します。Prolog はバッククォートの節と問い合わせ、
        Lean は宣言（\texttt{theorem} など）、KDB は q のコード、Filter はプロジェクト内のパスです。
  \item MATLAB と KDB では、日本語の文章をコードとして送ることはせず、呼び出しを失敗させます。
\end{itemize}

%======================================================================
\section{v0.1 で使えないもの}

次の構文は v0.1 に含まれず、構文エラーとして拒否されます。黙って受理されることはありません。
\begin{itemize}
  \item \texttt{case}、\texttt{await}、\texttt{type} 宣言と状態機械、\texttt{retry}、\texttt{approve}、
        \texttt{policy}、\texttt{export}、\texttt{configuration}、時相アサーション
  \item Event 宣言と完了条件（\texttt{completion when hard \{\dots\}}）の構文。意味論は実装済みで、
        完了条件は現在 \texttt{config/routing-defaults.toml} で与えます。
  \item 機密区分 \texttt{confidential}／\texttt{secret}、配置 \texttt{sandbox}／\texttt{private\_cloud}／\texttt{external}（E401）
\end{itemize}

%======================================================================
\appendix
\section{文法の早見表}

\begin{lstlisting}[language=AWHDL]
entity name is
    port ( a : in text<internal>; b : out data<public> );
end name;

architecture arch of name is
    device d : agent | mcp | deterministic
        generic (route => "<route>", location => local | cloud,
                 clearance => public | internal | restricted);
    signal s1, s2 : data<internal> := 0;
    timer t : period 1 sec;
    budget lim is iterations <= 10; wall_time <= 5 min;
                  tool_calls <= 3; model_calls <= 2; end budget;
    barrier bar (s1, s2);
begin
    assert always (<expr>);
    assert never (<expr>);
    assert never (internal -> cloud);

    process(a, s1.changed, d.done, d.failed, d.timeout, t, bar.ready)
        timeout 2 min;
    begin
        d.method(<expr>) timeout 30 sec -> s1;
        s2 <= s1 + 1;
        if <expr> then null; elsif <expr> then null; else null; end if;
        parallel d.m1(a) -> s1; d.m2(a) -> s2; end parallel;
        assert <expr>;
    on timeout
        b <= "timed out";
    end process;
end arch;
\end{lstlisting}

式の結合は弱い順に \texttt{or}、\texttt{and}、\texttt{not}、比較（\texttt{=}、\texttt{/=}、\texttt{<}、
\texttt{<=}、\texttt{>}、\texttt{>=}）、整数の \texttt{+}／\texttt{-} です。
比較は数値どうし、または文字列どうしで行います。

\section{診断コード}

\begin{longtable}{lp{108mm}}
\toprule
コード & 意味 \\
\midrule
\endhead
E001 & ソースファイルを読み取れない \\
E101 & 構文エラー \\
E201--E205 & 重複した entity・ポート・device・宣言、未知の entity \\
E206 & 感度リストの未知の名前 \\
E207／E208／E209 & 未知の device、未知の出力先 \\
E210 & 式の中の未知の名前 \\
E211 & 正でない時間（タイマー周期、制限時間） \\
E212 & 不正または重複した予算項目 \\
E214 & signal でないバリアのメンバー \\
E215 & 2 つの parallel 分岐が同じ signal に書き込む \\
E217 & 入力ポートへの書き込み \\
E218 & device の generic の重複や型の誤り \\
E301 & クラウドの device へ流せないデータ \\
E302 & device の clearance を超えるデータ \\
E303 & 高い区分のデータを低い区分の signal に保存 \\
E401 & v0.1 の Profile に含まれない構成要素 \\
E402 & 未知の名前（区分、配置など） \\
\bottomrule
\end{longtable}

\section{参考文献}

\begin{itemize}
  \item \texttt{awhdl/docs/PROFILE\_v0.1\_ja.md}：v0.1 の範囲と実装状況（正本）
  \item \texttt{awhdl/docs/LANGUAGE\_SPEC\_ja.md}、\texttt{AWHDL\_LANGUAGE\_SPEC\_v0.2\_ja.md}：言語仕様
  \item \texttt{awhdl/docs/RUNTIME\_SPEC\_ja.md}：delta cycle、Effect、完了評価
  \item \texttt{awhdl/docs/SECURITY\_SPEC\_ja.md}：機密区分、capability、承認
  \item \texttt{awhdl/docs/GETTING\_STARTED\_ja.md}：最初の一歩
\end{itemize}

\end{document}
